\documentclass[11pt,a4paper]{article}

\usepackage[margin=2.5cm]{geometry}
\usepackage[T1]{fontenc}
\usepackage{lmodern}
\usepackage{booktabs}
\usepackage{tabularx}
\usepackage{enumitem}
\usepackage{xcolor}
\usepackage{listings}
\usepackage{hyperref}

\hypersetup{colorlinks=true, linkcolor=blue!50!black, urlcolor=blue!50!black}
\lstset{
  basicstyle=\ttfamily\small,
  columns=fullflexible,
  breaklines=true,
  frame=single,
  framerule=0.3pt,
  rulecolor=\color{black!30},
  backgroundcolor=\color{black!3},
  keywordstyle=\color{blue!60!black}\bfseries,
  commentstyle=\color{black!55}\itshape,
  showstringspaces=false
}
\newcommand{\code}[1]{\texttt{#1}}

\title{CPU--GPU Co-Scheduling Module for CloudSim 7G\\[4pt]
\large Phase 0 and Phase 1a Report: Build Setup, Core Types, GPU-Layer Contract and Stub}
\author{Chirag}
\date{21 November 2026}

\begin{document}
\maketitle

\begin{abstract}
This report records the first implementation step of the co-scheduling module described in the
architecture document (frozen as v1). Phase~0 adds two Maven modules to CloudSim 7G and confirms
the whole project builds. Phase~1a writes the core types that every scheduling policy uses, the
contract between the scheduler and the GPU layer, and a stand-in (``stub'') GPU layer built from
the Check~1 prototypes. All 201 tests in the project pass, 21 of them new. While writing the stub
we found and fixed a systematic 0.01\,s start delay caused by CloudSim's event timing.
\end{abstract}

\tableofcontents

%======================================================================
\section{Context}
%======================================================================

The module lets a researcher describe a GPU cluster, a workload and candidate schedulers, and
compare how the schedulers perform, including how many CPU cores each training job receives.
Synergy (OSDI~'22) is replicated first and serves as the correctness check.

The architecture has five layers:

\begin{description}[leftmargin=2.2cm,style=nextline]
  \item[L1] GPU layer: GPU servers, guests, training jobs and their progress.
  \item[L2] Performance models: throughput, locality, interference seam, power.
  \item[L3] Workload: job traces, arrival generators, model labellers.
  \item[L4] Scheduling pipeline: ordering, allocation, placement, lifecycle, broker.
  \item[L5] Observation: metrics, energy, comparison runner, decision report.
\end{description}

L1 was originally assigned to a teammate. Because it had not been started, this phase defined
the L1 contract and a stub that satisfies it, so the scheduler can be developed and tested now.
As of 21~November~2026 the author owns L1 as well; the real GPU layer is the next step and will
replace the stub behind the same contract.

%======================================================================
\section{Phase 0: Build setup}
%======================================================================

\subsection{Repository}
The work lives in the existing CloudSim 7G clone (\code{cloudsim/}, upstream
\code{cloudslab/cloudsim}, version 7.0.1, Java~21) on the branch \code{cosched/skeleton}.
CloudSim's own code is not modified. The prototype programs \code{CheckA}--\code{CheckE.java}
from Check~1 are kept in \code{modules/} as a historical record.

\subsection{Maven modules}
Two modules were added to the root \code{pom.xml}:

\begin{center}
\begin{tabularx}{\linewidth}{@{}lX@{}}
\toprule
Module & Contents \\
\midrule
\code{cloudsim-gpu} & L1: the contract package \code{gpu.api} and the stand-in
  \code{gpu.stub}. Depends only on CloudSim. \\
\code{cloudsim-cosched} & Core, L2--L5. Depends on CloudSim and \code{cloudsim-gpu}. \\
\bottomrule
\end{tabularx}
\end{center}

The dependency runs in one direction only: the GPU module knows nothing about the scheduler.
This is why the contract uses its own small value types (Section~\ref{sec:api}).

\subsection{Dependencies}
Added to the root \code{dependencyManagement}:
\begin{itemize}
  \item \code{jackson-databind} and \code{jackson-dataformat-yaml} 2.17.2: reading the experiment
        YAML and the profile JSON.
  \item \code{archunit-junit5} 1.3.0 (test scope): enforcing the architecture rules in later phases.
  \item \code{opencsv} was already present and will read the job CSV.
\end{itemize}
Both new modules use Surefire 3.2.5 so JUnit~5 tests run.

\subsection{Package layout}
\begin{lstlisting}
org.cloudbus.cloudsim.gpu.api          contract with the GPU layer
org.cloudbus.cloudsim.gpu.stub         stand-in GPU layer
org.cloudbus.cloudsim.cosched.core     core types (this phase)
org.cloudbus.cloudsim.cosched.model    L2 (empty, next phases)
org.cloudbus.cloudsim.cosched.workload L3 (empty)
org.cloudbus.cloudsim.cosched.policy   L4 plug points (empty)
org.cloudbus.cloudsim.cosched.broker   L4 broker (empty)
org.cloudbus.cloudsim.cosched.metrics  L5 (empty)
\end{lstlisting}

%======================================================================
\section{Phase 1a: Core types}
%======================================================================

All classes are in \code{org.cloudbus.cloudsim.cosched.core}. Value types are Java records:
immutable, validated in their constructors, with errors that name the offending item.

\subsection{Resources and clusters}

\paragraph{\code{ResourceVector}} The amount a job holds: whole GPUs (integer), CPU cores and
memory in GB (doubles, because Synergy allocates fractional shares). It provides \code{plus},
\code{minus} (returning a \code{Delta} that may be negative), \code{fitsIn} with a tolerance of
$10^{-9}$ for floating-point noise, and \code{cpusPerGpu}.

\paragraph{\code{ServerGroup}} One line of the cluster YAML: name, count, GPUs, GPU type,
CPUs, memory, storage and network bandwidth. It also gives the GPU-proportional (``fixed
$m{:}n$'') share of a job:
\[
  \text{share}(g) = \Big(g,\ \tfrac{\text{cpus}}{\text{gpus}}\,g,\ \tfrac{\text{mem}}{\text{gpus}}\,g\Big),
\]
so a 2-GPU job on a Synergy server (8 GPUs, 24 CPUs, 500\,GB) gets 6 cores and 125\,GB.

\paragraph{\code{ClusterSpec}} The list of server groups plus an optional rack count. It expands
the groups into numbered \code{Server}s (group by group, in the order written), gives the total
capacity $(G, C, M)$ used by the DRF and Tetris keys, and lists the GPU types present so a later
check can stop the run when a profile is missing. \code{ClusterSpec.synergyDefault()} is Synergy's
cluster (16 servers $\times$ 8 V100, 24 CPUs, 500\,GB); \code{synergy(n)} varies the CPU count for
Fig.~12. The YAML example from the architecture document parses unchanged:

\begin{lstlisting}
server_groups:
  - { name: training-a100, count: 8,  gpus: 8, gpu_type: A100, cpus: 64,
      mem_gb: 1024, storage_mbps: 2000, net_gbps: 100 }
  - { name: legacy-v100,   count: 16, gpus: 8, gpu_type: V100, cpus: 24,
      mem_gb: 500,  storage_mbps: 500,  net_gbps: 40 }
racks: 2
\end{lstlisting}

\paragraph{\code{Allocation}} Where a job runs: one \code{Part} per server, each holding at least
one GPU. Two parts on the same server are rejected (they must be merged). \code{samePlacement}
tells whether two allocations use the same servers with the same GPUs, differing only in cores
or memory; the broker uses it to choose between an in-place resize and a move.

\paragraph{\code{Organisation}} Teams with optional GPU quotas, and priority classes. Without
teams it defaults to one team, \code{default}, with no quota.

\subsection{What policies see: \code{JobInfo} and \code{ClusterState}}
Architecture rule~6 says policies read only the cluster view and the model interfaces, and write
only to the round plan. Two read-only interfaces carry that view:

\begin{itemize}
  \item \code{JobInfo}: id, arrival time, GPUs requested, model, team, priority, total and
        remaining work, remaining time at the proportional share (for SRTF and Tetris), attained
        GPU-seconds (for LAS), first start time and current allocation.
  \item \code{ClusterState}: current time, servers, free resources per server, active jobs,
        organisation, and the busy GPU fraction (for Philly-style quota preemption at 90\%).
\end{itemize}

\subsection{\code{RoundPlan}}
A round plan is the draft of one scheduling round. It is the answer to the first risk found while
reading Synergy's code: Tune shrinks other jobs while placing one, which a read-only view cannot
allow. In the plan:

\begin{enumerate}
  \item Every \code{assign(job, allocation)} is checked against each server's free capacity
        at once. A write that does not fit changes nothing and returns \code{false}, so the plan
        is always feasible.
  \item A job may be assigned again in the same round; its own previous entry counts as free.
        This is how Tune shrinks an over-share job and then places the new one.
  \item Gang allocation is enforced: an allocation must hold exactly the GPUs the job asked for.
  \item A plan starts either \code{EMPTY} (Synergy's round-based mode: every running job must be
        re-allocated or it is preempted) or \code{CARRY\_RUNNING} (non-preemptive mode: running
        jobs keep their allocation and the round only adds waiting jobs).
\end{enumerate}

Free capacity is tracked as one used-resources vector per server, updated on every write, so
checks cost $O(\text{parts})$.

\subsection{\code{PolicyRegistry} and \code{Params}}
Policies are found by name (rule~7). The registry maps (kind, name) to a factory; an unknown name
fails with the list of known names. Each factory reads its parameters through \code{Params},
where every read states a default and a valid range:

\begin{lstlisting}[language=Java]
double roundS = p.getDouble("round_s", 300, 1, 86400);
\end{lstlisting}

This gives the architecture document's parameter rules directly. Values are range-checked at
startup; every resolved value, defaults included, is kept for the \code{RunManifest}; and after
construction any key the policy did not read is reported, which catches typos such as
\code{round\_sec}.

%======================================================================
\section{Phase 1a: The GPU-layer contract (\code{gpu.api})}\label{sec:api}
%======================================================================

\subsection{Design}
The scheduler decides; the GPU layer carries out. Each round, the broker turns the
\code{RoundPlan} into a list of commands for the GPU layer:

\begin{center}
\begin{tabularx}{\linewidth}{@{}lX@{}}
\toprule
Command & Meaning \\
\midrule
\code{Start(job, remainingWork, parts, rate)} & First start or resume of a waiting job. \\
\code{Resize(job, parts, rate)} & Change cores, memory or speed. Same servers and GPUs: resized
  in place. Otherwise: the job is moved (stopped and started with its remaining work). \\
\code{Stop(job)} & Preempt: release resources and keep the remaining work. \\
\bottomrule
\end{tabularx}
\end{center}

\emph{Rate} is the job's speed in iterations per second. It already includes throughput, locality
and any interference factor. The GPU layer turns rate into CloudSim speed and never computes
throughput itself. This is the ``rate-to-speed split'' agreed in Check~3.

\paragraph{Why separate value types.} The architecture document lists \code{ResourceVector} as
shared with L1. Using it in \code{gpu.api} would make the GPU module depend on the scheduler
module, which depends on the GPU module: a cycle. The contract therefore uses \code{PartSpec}
(server, GPUs, cores, memory) and \code{ServerShape} (one server to build), and the broker maps
between them. This is the only change from the architecture document in this phase.

\subsection{Interfaces}
\begin{itemize}
  \item \code{GpuLayer}: \code{build(servers)} creates the datacenter and one host per server;
        \code{bind(broker)} returns the executor that acts for that broker.
  \item \code{PlanExecutor}: \code{apply(commands)}, \code{remainingWork(job)},
        \code{handle(event)} (the broker forwards every event; the GPU layer consumes its own),
        and \code{setJobEndListener}.
  \item \code{JobEndListener}: called once per job when all its parts have finished.
\end{itemize}

\subsection{Rules every implementation must follow}
These come from Check~1 and are written into the \code{PlanExecutor} documentation:

\begin{enumerate}
  \item Before any resize or removal, credit every running job's progress at its old speed.
  \item A resize is deallocate, then allocate; allocating twice leaks host capacity.
  \item After a resize, wake the datacenter so finish times are recomputed.
  \item Placement is exactly what the command says; a full host refuses, nothing falls back.
  \item A refused creation or resize is a bug and stops the run.
  \item The executor never ends the simulation; the broker does.
\end{enumerate}

Within one \code{apply}, all stops run first, then resizes, then starts, so capacity freed in a
round can be reused in the same round regardless of the order the commands are listed in.

%======================================================================
\section{Phase 1a: The stub GPU layer (\code{gpu.stub})}
%======================================================================

The stub is the Check~1 prototype code (mainly \code{CheckE}) turned into classes behind
the contract. It exists so the scheduler can run end to end before the real GPU layer is ready.

\begin{itemize}
  \item \code{StubGpuHost} extends CloudSim's \code{PowerHost}. CloudSim sees one processing
        element of $10^{15}$\,MIPS, large enough never to limit speed; GPUs, cores and memory
        are counted by the host itself, and any over-commit throws (rule~5).
  \item \code{PlannedGuestAllocationPolicy} places each guest exactly on the host the plan names
        or refuses it (rule~4). In Check~1, CloudSim's default policy put 15 of 20 guests
        off plan; this one puts none.
  \item \code{StubPlanExecutor}: each job part is one VM with one cloudlet holding the job's full
        remaining work. All parts run at the job's rate, so a spread job's parts finish together
        (lock-step). Speed is carried as VM MIPS with 1 iteration $= 10^6$\,MI, which keeps
        sub-iteration precision in CloudSim's integer cloudlet lengths.
\end{itemize}

\subsection{Bug found: every job started 0.01\,s late}
The first test run showed a job's remaining work at 100\,s as 900.01 instead of 900 (rate 1, work
1000). The cause is in CloudSim's \code{Datacenter} (line~413): the VM-creation acknowledgement is
sent after \code{minTimeBetweenEvents}, 0.01\,s by default. The Check~1 prototypes, like the first
version of the stub, submitted the cloudlet only after that acknowledgement, so every job began
0.01\,s after its start command. The Check~1 tolerances (0.05\,s) were loose enough to hide it;
the stub's tests check remaining work directly, which exposed it.

\textbf{Fix.} The cloudlet is now submitted immediately after the VM-creation request. The
datacenter processes both at the same simulated instant, in order, so the VM exists when the
cloudlet arrives. The acknowledgement is used only to detect a refused creation. After the fix,
remaining work and finish times match the hand-computed values exactly.

%======================================================================
\section{Tests}
%======================================================================

\subsection{Stub GPU layer (8 tests)}
A small driver entity sends scripted commands at fixed times and records job ends.

\begin{center}
\begin{tabularx}{\linewidth}{@{}lX@{}}
\toprule
Test & Checks \\
\midrule
constant rate & 1000 iterations at 2/s finish at 500\,s; resources released. \\
in-place resize & Rate 1, then rate 4 and 18 cores at 100\,s: remaining work 900 at 100\,s,
  finish at $100 + 900/4 = 325$\,s, host cores updated. \\
stop and resume & Stopped at 100\,s with 900 left; resumed on another server at 200\,s at
  rate 3; finishes at 500\,s. \\
move & Resize to another server releases the old one and continues with the remaining work. \\
same-round reuse & A stop and a start in one round, listed start-first, still fit. \\
over-commit & Two jobs needing 9 GPUs on an 8-GPU server stop the run with an error. \\
spread job & Two parts on two servers finish together; the job ends exactly once. \\
random resizes & CheckE ported: 300 jobs with 1, 2 or 4 parts, a new random rate every 300\,s
  round, finish times compared with an independent calculation; worst error below 0.05\,s. \\
\bottomrule
\end{tabularx}
\end{center}

\subsection{Core types (12 tests)}
Cluster expansion and total capacity; parsing the architecture document's cluster YAML; validation
messages naming the bad server group; the proportional share; \code{RoundPlan} rejecting what does
not fit and enforcing gangs; revision within a round (the Tune pattern); \code{CARRY\_RUNNING}
versus \code{EMPTY}; allocations with two parts on one server rejected; \code{samePlacement};
registry lookup and known-name errors; parameter range checks and typo detection; organisation
defaults and duplicate teams.

\subsection{Whole project}
\begin{center}
\begin{tabular}{@{}lrr@{}}
\toprule
Module & Tests & Failures \\
\midrule
\code{cloudsim} (core, unchanged) & 167 & 0 \\
\code{cloudsim-examples} & 13 & 0 \\
\code{cloudsim-gpu} & 8 & 0 \\
\code{cloudsim-cosched} & 13 & 0 \\
\midrule
Total & 201 & 0 \\
\bottomrule
\end{tabular}
\end{center}
Command: \code{mvn install} from \code{cloudsim/}.

%======================================================================
\section{Next steps}
%======================================================================
\begin{enumerate}
  \item \textbf{Real GPU layer} (now owned by the author): \code{GpuHost} (PowerHost with GPU
        provisioning, GPU type and power), \code{GpuProvisioner} (gang allocation), \code{GpuVm},
        \code{TrainingJob} and \code{TrainingCloudlet}, \code{GpuCloudletScheduler} and a
        \code{GpuPlanExecutor}, replacing the stub behind \code{gpu.api}. The stub's tests become
        the acceptance tests for the real layer.
  \item \textbf{Formats and loaders}: job CSV, profile JSON, experiment YAML, documented in
        \code{docs/formats.md}.
  \item \textbf{L2 models}: \code{ThroughputModel}, \code{ProfileRepository}, \code{LocalityModel};
        the Synergy profile converter and the fidelity test (ResNet-18 at 9 cores $= 2.4\times$),
        which need Synergy's repository in \code{Experiment/synergy/}.
\end{enumerate}

\end{document}
